\section{Round-7 Objective}
We pursue \textbf{(A) a proof strategy}.  Round~6 produced a sharp local-clique criterion for existence of
\(\delta_{\log}(A)\), but that argument still depended on the compatibility graph and on counting compatible
subfamilies.

The new goal of this round is to remove \emph{all} compatibility bookkeeping in a substantial regime.
The key observation is that for each fixed stage \(k\), the true set \(A\) is squeezed below the periodic stage-set
\(A^{(k)}\), and the only possible loss comes from \emph{individual later congruence classes}.
A crude union bound therefore yields a new theorem depending only on the tail of
\(\sum 1/n_i\), with no reference to clique counts or lcm-growth.

The main consequence is:
\begin{equation}
\sum_{i=1}^{\infty}\frac{1}{n_i}<\infty
\qquad\Longrightarrow\qquad
\delta_{\log}(A)\ \text{exists}.
\tag{7.1}
\end{equation}
This is a genuine advance beyond Round~6, because it applies even when the compatibility graph is complete and the
Round~6 clique counts are exponential.

\section{Round-6 Foundation Used}
We rely on the following earlier results.
\begin{itemize}
\item \textbf{Activation identity} (Round~1): for
\[
A^{(k)}:=\{n\in\mathbb N:\ n\not\equiv a_i\pmod{n_i}\ \forall\,1\le i\le k\},
\]
one has
\[
\mathbf 1_A(n)=\mathbf 1_{A^{(k)}}(n)\qquad (n_k\le n<n_{k+1}).
\]
\item \textbf{Periodic stage densities} (Round~1 and Round~2): each \(A^{(k)}\) is periodic, hence has a natural density
\[
\delta_k:=d(A^{(k)}),
\]
and the sequence \((\delta_k)\) is decreasing with a limit
\[
\delta:=\lim_{k\to\infty}\delta_k\in[0,1].
\]
\item \textbf{Logarithmic asymptotic for a periodic stage-set} (already used repeatedly in earlier rounds):
for fixed \(k\),
\[
S_k(x):=\sum_{\substack{1\le n\le x\\ n\in A^{(k)}}}\frac1n
=\delta_k\log x+O_k(1),
\]
so
\[
\frac{S_k(x)}{\log x}\to \delta_k.
\]
\item \textbf{Single progression harmonic estimate} (Round~5): for every modulus \(m\ge 1\), every residue class
\(b\pmod m\), and integers \(2\le X<Y\),
\[
\sum_{\substack{X\le n<Y\\ n\equiv b\ (\mathrm{mod}\ m)}}\frac1n
=
\frac1m\log\frac{Y}{X}+O\!\left(\frac1X\right),
\]
with an absolute implied constant, uniform in \(m,b,X,Y\).
\end{itemize}

\section{New Insight / Tool (Round-7)}
The new tool is a \emph{fixed-stage sandwich}.
Instead of expanding \(A\) globally by inclusion--exclusion (Round~6), we compare \(A\) directly to a fixed finite stage
\(A^{(k)}\).
Since
\[
A\subseteq A^{(k)},
\]
any integer counted by \(A^{(k)}\) but not by \(A\) must violate \emph{some later individual constraint}
\(n\equiv a_i\pmod{n_i}\) with \(i>k\).
Thus the entire error is bounded by a union of single arithmetic progressions.

This has two major consequences.
\begin{itemize}
\item The bound is completely independent of the compatibility graph and of all inclusion--exclusion clique counts.
\item Any counterexample must satisfy
\begin{equation}
\sum_{i=1}^{\infty}\frac1{n_i}=\infty,
\tag{7.2}
\end{equation}
since the convergent-reciprocal regime is now ruled out unconditionally.
\end{itemize}

\section{Attack Plan (Round-7)}
The exact gap left after Round~6 was this: even after charging each compatible subfamily only once, the error term still
involved local clique counts \(C_r\), which can be enormous when many congruences are mutually compatible.

The Round~7 plan is:
\begin{enumerate}
\item Fix \(k\) and compare \(A\) with the periodic set \(A^{(k)}\).
\item Prove a one-sided pointwise inequality showing that every integer in \(A^{(k)}\setminus A\) is caught by at least
one later single progression.
\item Sum with weight \(1/n\), use the Round~5 single-progression estimate, and derive a general tail criterion.
\item Deduce the clean corollary (7.1) when \(\sum 1/n_i\) converges.
\item Give an explicit family showing that this new theorem applies far beyond the Round~6 clique criterion.
\end{enumerate}

\section{Work (Round-7)}
Throughout this round, if \(n_1=1\) then \(A=\varnothing\), so \(\delta_{\log}(A)=0\) trivially.  Hence we may assume
\begin{equation}
n_1\ge 2.
\tag{7.3}
\end{equation}

For \(x\ge 1\), write
\begin{equation}
S(x):=\sum_{\substack{1\le n\le x\\ n\in A}}\frac1n,
\qquad
S_k(x):=\sum_{\substack{1\le n\le x\\ n\in A^{(k)}}}\frac1n.
\tag{7.4}
\end{equation}

\paragraph{Lemma 7.1 (fixed-stage domination by later single constraints).}
For every fixed \(k\ge 0\) and every integer \(n\ge 1\),
\begin{equation}
0\le \mathbf 1_{A^{(k)}}(n)-\mathbf 1_A(n)
\le
\sum_{i>k}\mathbf 1_{\,n\ge n_i\,}\,\mathbf 1_{\,n\equiv a_i\ (\mathrm{mod}\ n_i)}.
\tag{7.5}
\end{equation}
Consequently, for every \(x\ge 1\),
\begin{equation}
0\le S_k(x)-S(x)
\le
\sum_{\substack{i>k\\ n_i\le x}}
\sum_{\substack{n_i\le n\le x\\ n\equiv a_i\ (\mathrm{mod}\ n_i)}}\frac1n.
\tag{7.6}
\end{equation}

\emph{Proof.}
The left inequality is immediate from \(A\subseteq A^{(k)}\).
If \(n\in A^{(k)}\setminus A\), then \(n\) avoids all congruences with indices \(\le k\), but fails at least one later active
constraint.  That is, there exists some index \(i>k\) such that \(n_i\le n\) and \(n\equiv a_i\pmod{n_i}\).
This proves (7.5).  Summing (7.5) for \(1\le n\le x\) with weight \(1/n\) gives (7.6).  \hfill
\(\square\)

\paragraph{Proposition 7.2 (general tail criterion).}
For each \(k\ge 0\), define
\begin{equation}
\Theta_k
:=
\limsup_{x\to\infty}
\frac1{\log x}
\sum_{\substack{i>k\\ n_i\le x}}
\frac{\log(x/n_i)+1}{n_i}.
\tag{7.7}
\end{equation}
If
\begin{equation}
\Theta_k\longrightarrow 0\qquad(k\to\infty),
\tag{7.8}
\end{equation}
then \(\delta_{\log}(A)\) exists and equals
\begin{equation}
\delta:=\lim_{k\to\infty}\delta_k.
\tag{7.9}
\end{equation}

\emph{Proof.}
Fix \(k\ge 0\).  By Lemma~7.1 and the Round~5 single-progression estimate,
\begin{align}
0\le S_k(x)-S(x)
&\le
\sum_{\substack{i>k\\ n_i\le x}}
\left(
\frac1{n_i}\log\frac{x}{n_i}+O\!\left(\frac1{n_i}\right)
\right) \notag\\
&\le
\sum_{\substack{i>k\\ n_i\le x}}
\frac{\log(x/n_i)+C}{n_i}
\tag{7.10}
\end{align}
for an absolute constant \(C\).  Replacing \(+1\) in (7.7) by \(+C\) clearly changes \(\Theta_k\) by at most an absolute
multiplicative constant, so (7.8) is equivalent with \(+C\) in place of \(+1\).
Therefore
\begin{equation}
\limsup_{x\to\infty}\frac{S_k(x)-S(x)}{\log x}\le \Theta_k.
\tag{7.11}
\end{equation}
Since \(S_k(x)=\delta_k\log x+O_k(1)\), we obtain
\begin{equation}
\delta_k-\Theta_k
\le
\liminf_{x\to\infty}\frac{S(x)}{\log x}
\le
\limsup_{x\to\infty}\frac{S(x)}{\log x}
\le
\delta_k.
\tag{7.12}
\end{equation}
Now let \(k\to\infty\).  Round~2 gives \(\delta_k\downarrow\delta\), and by hypothesis \(\Theta_k\to 0\).  Hence
both the liminf and limsup of \(S(x)/\log x\) are squeezed to \(\delta\).  Thus
\[
\frac{S(x)}{\log x}\to \delta,
\]
which is exactly \(\delta_{\log}(A)=\delta\).  \hfill \(\square\)

\paragraph{Theorem 7.3 (summable reciprocal moduli imply existence).}
If
\begin{equation}
\sum_{i=1}^{\infty}\frac1{n_i}<\infty,
\tag{7.13}
\end{equation}
then \(\delta_{\log}(A)\) exists and equals \(\delta=\lim_k\delta_k\).
More precisely, for every fixed \(k\ge 0\),
\begin{equation}
\delta_k-\sum_{i>k}\frac1{n_i}
\le
\liminf_{x\to\infty}\frac{S(x)}{\log x}
\le
\limsup_{x\to\infty}\frac{S(x)}{\log x}
\le
\delta_k.
\tag{7.14}
\end{equation}

\emph{Proof.}
Under (7.13), the tail \(\sum_{i>k}1/n_i\) is finite.  For every \(x\),
\begin{align}
\frac1{\log x}
\sum_{\substack{i>k\\ n_i\le x}}
\frac{\log(x/n_i)+1}{n_i}
&\le
\sum_{i>k}\frac1{n_i}
+
\frac1{\log x}\sum_{i>k}\frac1{n_i}.
\tag{7.15}
\end{align}
Taking \(x\to\infty\) yields
\begin{equation}
\Theta_k\le \sum_{i>k}\frac1{n_i}.
\tag{7.16}
\end{equation}
Hence \(\Theta_k\to 0\) as \(k\to\infty\), so Proposition~7.2 applies and gives existence.
Combining (7.12) with (7.16) gives the sharper bound (7.14).  \hfill \(\square\)

\paragraph{Corollary 7.4 (a large new class).}
If
\begin{equation}
n_i\gg i^{1+\eta}\qquad\text{for some }\eta>0,
\tag{7.17}
\end{equation}
then \(\delta_{\log}(A)\) exists for \emph{every} choice of residues \(a_i\pmod{n_i}\).
In particular, every sequence with at least polynomial growth of exponent strictly greater than \(1\) is excluded as a
possible counterexample.

\emph{Proof.}
Condition (7.17) implies \(\sum_i 1/n_i<\infty\), so Theorem~7.3 applies.  \hfill \(\square\)

\paragraph{Example 7.5 (a family far beyond the Round~6 clique criterion).}
Let \(p_i\) be the \(i\)-th prime, and set
\begin{equation}
n_i:=p_i^2\qquad(i\ge 1),
\tag{7.18}
\end{equation}
with arbitrary residues \(a_i\pmod{p_i^2}\).
Then \(\delta_{\log}(A)\) exists by Theorem~7.3, because
\[
\sum_{i=1}^{\infty}\frac{1}{p_i^2}<\infty.
\]
However, this family is \emph{not} covered by the earlier structural criteria.

Indeed, the moduli \(p_i^2\) are pairwise coprime, so every finite subfamily of congruences is compatible.  Thus the
compatibility graph is complete, and the local clique counts from Round~6 satisfy
\begin{equation}
C_r=2^{r-1}.
\tag{7.19}
\end{equation}
Hence the Round~6 local-count condition fails badly:
\[
\sum_{r\le K}\frac{C_r}{n_r}
\ge
\sum_{r\le K}\frac{2^{r-1}}{p_r^2},
\]
and the summands themselves grow exponentially relative to \(p_r^2\sim r^2(\log r)^2\).
Also the Round~4 bounded-lcm hypothesis fails, since
\[
\frac{\mathrm{lcm}(n_1,\dots,n_k)}{n_k}
=
\prod_{i<k}p_i^2
\to\infty.
\]
Thus Theorem~7.3 genuinely covers a new regime beyond all previous rounds.

\section{Adversarial Verification}
We stress-test the new argument.
\begin{itemize}
\item \textbf{One-sided union bound only.}  The proof does not claim equality in (7.5); multiple later violations may
occur simultaneously.  This is harmless because only an upper bound for \(S_k(x)-S(x)\) is needed.
\item \textbf{No compatibility assumptions whatsoever.}  Lemma~7.1 uses only the implication
\(n\in A^{(k)}\setminus A\Rightarrow\exists i>k\) with \(n\equiv a_i\pmod{n_i}\).  No lcm control, no clique counts, and no
Chinese remainder arguments are required.
\item \textbf{Use of the single-progression estimate.}  The Round~5 estimate is uniform in the modulus and residue class, so it
is legitimate to sum it over all late indices \(i>k\).
\item \textbf{Stage asymptotic.}  The passage \(S_k(x)=\delta_k\log x+O_k(1)\) is valid because \(A^{(k)}\) is periodic; this was
already established in earlier rounds.
\item \textbf{Edge case \(n_1=1\).}  Treated separately at the start: then \(A=\varnothing\), so the theorem is trivially true.
\item \textbf{Strict advance beyond Round~6.}  Example~7.5 shows that Theorem~7.3 applies in a family where the Round~6 local
counts are exponential and the bounded-lcm framework from Round~4 also fails.
\item \textbf{What remains open.}  The new theorem leaves untouched the genuinely difficult regime
\(\sum 1/n_i=\infty\).  This is consistent with the proof, since the union bound only gives useful control when the tail of
later single progressions is harmonically small.
\end{itemize}

\section{Final}
\textbf{UNRESOLVED (BUT STRICTLY ADVANCED).}

The key new result is Theorem~7.3:
\begin{quote}
If \(\sum_i 1/n_i<\infty\), then the logarithmic density \(\delta_{\log}(A)\) exists for every choice of residues.
\end{quote}
This is a substantial improvement over Round~6 because it eliminates all compatibility-graph and clique-count hypotheses.
In particular, it rules out an entire class of counterexample strategies:
\begin{equation}
\text{any counterexample must satisfy }\sum_{i=1}^{\infty}\frac1{n_i}=\infty.
\tag{7.20}
\end{equation}
So after Round~7, the unresolved core has narrowed further:
counterexamples, if they exist at all, must occur in the divergent-reciprocal regime, where neither the Round~4 lcm
control, nor the Round~6 local-clique control, nor the new Round~7 tail-union argument applies.

\section{Completion Estimate}
\textbf{COMPLETION: 91\%}.

\section{References}
\begin{itemize}
\item No external references were used in this round.
\item The argument builds directly on the vetted Round~1--Round~6 results supplied in the conversation.
\end{itemize}
